\subsection{Projectieve representaties}
De \(U(g)\)'s vermenigvuldigen elkaar slechts tot op een fase, en vormen dus een projectieve representatie van de groep \(G\).
 Dit betekent dat voor twee groepselementen \(g\) en \(h\) de operatoren \(U(g)\) en \(U(h)\) vermenigvuldigen volgens:
    \[
        U(g)U(h) = e^{i\omega(g,h)}U(gh)
    \]

Een nodige conditie is \textbf{associativiteit}, wat inhoudt dat:
    \[
        \omega(g,h) + \omega(gh,k) = \omega(h,k) + \omega(g,hk)
    \]

\begin{derivation}
    \[(U(g)U(h))U(k) = U(g)(U(h)U(k)) \]

\begin{align*}
   (U(g)U(h))U(k) &= e^{i\omega(g,h)}U(gh)U(k) \\
    &= e^{i\omega(g,h)}e^{i\omega(gh,k)}U(ghk) \\
\end{align*}
\begin{align*}
    U(g)(U(h)U(k)) &= U(g)e^{i\omega(h,k)}U(hk) \\
    &= e^{i\omega(h,k)}e^{i\omega(g,hk)}U(ghk)
\end{align*}
\[
\implies e^{i\omega(g,h)}e^{i\omega(gh,k)} = e^{i\omega(h,k)}e^{i\omega(g,hk)}
\]
\[
    \implies \omega(g,h) + \omega(gh,k) = \omega(h,k) + \omega(g,hk)
\]
\end{derivation}